
\documentclass{article}
\usepackage{colortbl}
\usepackage{makecell}
\usepackage{multirow}
\usepackage{supertabular}

\begin{document}

\newcounter{utterance}

\centering \large Interaction Transcript for game `bbh', experiment `sports\_understanding', episode 67 with reflexion{-}Qwen3.5{-}4B.
\vspace{24pt}

{ \footnotesize  \setcounter{utterance}{1}
\setlength{\tabcolsep}{0pt}
\begin{supertabular}{c@{$\;$}|p{.15\linewidth}@{}p{.15\linewidth}p{.15\linewidth}p{.15\linewidth}p{.15\linewidth}p{.15\linewidth}}
    \# & \multicolumn{2}{c}{Player} && \multicolumn{2}{c}{Game Master} \\
    \hline

    \theutterance \stepcounter{utterance}  
    & & & \multicolumn{4}{p{0.6\linewidth}}{
        \cellcolor[rgb]{0.9,0.9,0.9}{
            \makecell[{{p{\linewidth}}}]{
                \texttt{\tiny{[P1$\langle$GM]}}
                \texttt{Determine whether an artificially constructed sentence relating to sports is plausible or not.} \\
\\ 
\texttt{Q: Is the following sentence plausible? "Bam Adebayo scored a reverse layup in the Western Conference Finals."} \\
\texttt{A: yes} \\
\\ 
\texttt{Q: Is the following sentence plausible? "Santi Cazorla scored a touchdown."} \\
\texttt{A: no} \\
\\ 
\texttt{Q: Is the following sentence plausible? "DeMar DeRozan was called for the goal tend."} \\
\texttt{A: yes} \\
\\ 
\texttt{Q: Is the following sentence plausible? "Jordan Binnington scored in the third period."} \\
\texttt{A:} \\
            }
        }
    }
     \\ \\

    \theutterance \stepcounter{utterance}  
    & \multicolumn{4}{p{0.6\linewidth}}{
        \cellcolor[rgb]{0.561,0.737,0.561}{
            \makecell[{{p{\linewidth}}}]{
                \texttt{\tiny{[P1$\rangle$GM]}}
                \texttt{no} \\
            }
        }
    }
    & & \\ \\

    \theutterance \stepcounter{utterance}  
    & & & \multicolumn{2}{p{0.3\linewidth}}{
        \cellcolor[rgb]{0.9,0.9,0.9}{
            \makecell[{{p{\linewidth}}}]{
                \texttt{\tiny{[GM$|$GM]}}
                \texttt{Target: yes} \\
            }
        }
    }
    & & \\ \\

    \theutterance \stepcounter{utterance}  
    & & & \multicolumn{2}{p{0.3\linewidth}}{
        \cellcolor[rgb]{0.9,0.9,0.9}{
            \makecell[{{p{\linewidth}}}]{
                \texttt{\tiny{[GM$|$GM]}}
                \texttt{game\_result = LOSE} \\
            }
        }
    }
    & & \\ \\

    \theutterance \stepcounter{utterance}  
    & & \multicolumn{4}{p{0.6\linewidth}}{
        \cellcolor[rgb]{0.95,0.95,0.95}{
            \makecell[{{p{\linewidth}}}]{
                \texttt{\tiny{[P1$\leftrightarrow$P]}}
                \texttt{Post{-}episode Reflexion (not sent to GM):} \\
\texttt{New plan: I failed to correctly identify the plausibility of the sentence because I did not verify the specific sport associated with the player and the scoring method. "Jordan Binnington" is a goaltender for the NHL, and the "third period" is a hockey term, not a soccer term (where he would be a goalkeeper). The sentence mixes hockey terminology with soccer context or implies a soccer event where he played, which is impossible. Next time, I will first identify the player's primary sport and the specific rules of that sport (e.g., scoring methods, game periods) before judging the sentence's plausibility.} \\
            }
        }
    }
    & \\ \\

\end{supertabular}
}

\end{document}
